\documentclass{article}
\usepackage[UTF8]{ctex}
\usepackage{amsmath,amssymb,amsfonts}
\usepackage{enumitem}
\usepackage{geometry}
\geometry{a4paper, margin=1in}
% === 错题本排版优化（不参与题目标识） ===
\setlength{\parindent}{0pt}
\setlength{\parskip}{0.35em plus 0.10em minus 0.08em}
\setlength{\emergencystretch}{3em}
\setlist{itemsep=0.25em, topsep=0.40em, parsep=0pt}
\allowdisplaybreaks[2]
% === 错题本排版优化结束 ===
\begin{document}
\section*{同济大学2025年数学分析真题}
\section*{40．同济大学2025年数学分析试卷}

\subsection{题目1}
1．（10 分）设 \(\alpha, \beta\) 为正实数，证明：

\[
\lim _{n \rightarrow \infty}\left(1+\frac{\alpha+\sqrt{-1} \beta}{n}\right)^{n}=e^{\alpha}(\cos \beta+\sqrt{-1} \sin \beta) .
\]

注：若数列 \(\left\{a_{n}\right\},\left\{b_{n}\right\}\) 分别收敛于 \(a, b\) ，则称复数列 \(\left\{a_{n}+\sqrt{-1} b_{n}\right\}\) 收敛于 \(a+\sqrt{-1} b\) ，记作

\[
\lim _{n \rightarrow \infty}\left(a_{n}+\sqrt{-1} b_{n}\right)=a+\sqrt{-1} b
\]

\subsection{题目2}
2．（10 分）证明：存在 \(f(x) \in C^{1}(\mathbb{R})\) ，满足 \(f(x)\) 在 \((-\infty, 1)\) 上恒为 0 ，在 \((2,+\infty)\) 上恒为 1 。

\begin{itemize}

\subsection{题目3}
\end{itemize}

\subsection{题目3}
原题边界“3．”：（15 分）设 \(f\) 在 \(\mathbb{R}\) 上二阶可微，证明：对任意的 \(a, b \in \mathbb{R}, t \in(0,1)\) ，存在 \(\xi \in \mathbb{R}\) ，满足
\[
f((1-t) a+t b)=(1-t) f(a)+t f(b)-\frac{f^{\prime \prime}(\xi)}{2} t(1-t)(b-a)^{2}
\]

\subsection{题目4}

\subsection{题目4}
原题边界“4．”：（15 分）设 \(\alpha\) 为正实数，讨论 \(\lim _{x \rightarrow 0} \frac{\frac{1}{\sqrt{\pi}} \int_{0}^{+\infty} e^{-t^{2}} \cos (2 x t) \mathrm{d} t-\frac{1}{2}}{x^{\alpha}}\) 的敛散性及极限．

\subsection{题目5}
5．（15 分）设 \(u_{n}(x) \in R[0,1]\) ，且 \(\sum_{n=1}^{\infty} u_{n}(x)\) 的部分和函数列一致有界．令 \(v_{n}(x)=\int_{0}^{1} u_{n}(x) \sqrt{x+t} \mathrm{~d} t\) ．

\begin{itemize}
  \item[（1）] 证明：\(v_{n}(x) \in C^{1}(0,1)\) ．
  \item[] \((2)\) 证明：若 \(\sum_{n=1}^{\infty} v_{n}(x)\) 收敛于 \(S(x)\) ，则 \(S(x) \in C^{0}(0,1)\) ．
\end{itemize}

\subsection{题目6}
6．（15 分）设 \(f, g \in C^{0}\left(\mathbb{R}^{2}\right)\) ，且当 \(x \neq 0\) 时，成立 \(f(x)=g(x) \sin \frac{1}{\|x\|}\) ．若 \(f\) 在原点可微，求 \(\mathrm{d} f(0)\) ，并分析 \(g\) 在原点的可微性．

\subsection{题目7}
7．（15 分）判断 \(\iint_{x+y \geq 1} \frac{\sin x \sin y}{(x+y)^{3}} \mathrm{~d} x \mathrm{~d} y\) 的敛散性。

\subsection{题目8}
8．（10 分）曲线 \(L\) 是 \(x^{2}+y^{2}+z^{2}=8\) 与 \(x^{2}+y^{2}=2 z\) 的交线，且从 \(z\) 轴正向往坐标原点看去取逆时针方向．计算

\[
I=\int_{L}|\sqrt{3} y-x| \mathrm{d} x+3 y \mathrm{~d} y-2 z \mathrm{~d} z
\]

\subsection{题目9}
9．（25 分）设 \(f(x)\) 是在 \(x_{0}\) 局部有定义的实函数，如果 \(f(x)\) 在 \(x_{0}\) 的某邻域上可以展开为收敛于 \(f(x)\)的关于 \(\left(x-x_{0}\right)\) 的幂级数，则称 \(f(x)\) 在 \(x_{0}\) 点解析，若 \(f(x)\) 在开区间 \(I\) 上点点解析，则称 \(f(x)\) 在 \(I\)上解析。\\
（1）若 \(f, g\) 都为区间 \(I\) 上的实解析函数，且 \(E=\{x \in I \mid f(x)=g(x)\}\) 在 \(I\) 上有聚点，证明：\(f \equiv g\) ．\\
（2）设 \(f, g\) 为 \(\mathbb{R}\) 上的实解析函数，\(w(f, g)=\operatorname{det}\left(\begin{array}{cc}f & g \\ f^{\prime} & g^{\prime}\end{array}\right)\) 为 Wronski 行列式，证明：\(w(f, g)=0\) 当且仅当 \(f, g\) 线性相关．\\
（3）若 \(f(x), g(x) \in C^{\infty}(\mathbb{R})\) ，（2）中结论是否成立？说明理由．

\subsection{题目10}
10．（20 分）设 \(P(x, y), Q(x, y)\) 为单连通区域 \(\Omega\) 上的 \(C^{1}\) 函数，且 \(\frac{\partial P}{\partial y}=\frac{\partial Q}{\partial x}\) 在 \(\Omega\) 上成立．证明：曲线积分 \(\int_{L} P \mathrm{~d} x+Q \mathrm{~d} y\) 与路径无关，其中 \(L\) 为 \(\Omega\) 中的分段光滑曲线。
\end{document}
